\documentclass[11pt,a4paper]{article}

\usepackage[utf8]{inputenc}
\usepackage[T1]{fontenc}
\usepackage[a4paper,margin=2.4cm]{geometry}
\usepackage{lmodern}
\usepackage{xcolor}
\usepackage{booktabs}
\usepackage{array}
\usepackage{tabularx}
\usepackage{amsmath}
\usepackage{listings}
\usepackage{hyperref}

\definecolor{navy}{HTML}{17365D}
\definecolor{lightblue}{HTML}{EAF2F8}
\definecolor{codegray}{HTML}{F5F5F5}

\hypersetup{
  colorlinks=true,
  linkcolor=navy,
  urlcolor=navy,
  pdftitle={Viabilidad de RAG para Dispute Transaction Support},
  pdfauthor={Factored Hackathon 2026}
}

\lstset{
  basicstyle=\ttfamily\small,
  backgroundcolor=\color{codegray},
  frame=single,
  rulecolor=\color{lightblue},
  columns=fullflexible,
  keepspaces=true,
  showstringspaces=false,
  breaklines=true,
  literate={á}{{\'a}}1 {é}{{\'e}}1 {í}{{\'i}}1 {ó}{{\'o}}1 {ú}{{\'u}}1 {ñ}{{\~n}}1 {Á}{{\'A}}1 {É}{{\'E}}1 {Í}{{\'I}}1 {Ó}{{\'O}}1 {Ú}{{\'U}}1 {Ñ}{{\~N}}1
}

\renewcommand{\arraystretch}{1.18}
\setlength{\parindent}{0pt}
\setlength{\parskip}{6pt}

\title{\textcolor{navy}{Viabilidad de RAG para\\[2pt] Dispute Transaction Support}}
\author{Factored Hackathon 2026}
\date{27 de septiembre de 2026}

\begin{document}

\maketitle

\begin{abstract}
Este reporte evalúa si los datos disponibles en BigQuery pueden sustentar una solución de pregunta--respuesta para Dispute Transaction Support antes de recurrir a inferencia estadística general. La conclusión es que existe una oportunidad clara para un sistema híbrido: BigQuery como fuente estructurada y autorizada para datos actuales de tarjetas, transacciones y disputas; RAG sobre documentación oficial; y recuperación de casos históricos como evidencia secundaria. Los transcripts y reclamos actuales no son suficientes para funcionar como una base de conocimiento textual autónoma.
\end{abstract}

\section{Conclusión ejecutiva}

Los datos sí pueden fortalecer una solución de pregunta--respuesta, pero no conviene convertir todas las tablas de BigQuery en documentos vectoriales. El uso adecuado depende del tipo de pregunta:

\begin{itemize}
  \item preguntas sobre el estado actual de una tarjeta, una transacción o un reclamo: consulta estructurada a BigQuery;
  \item preguntas sobre políticas, requisitos y procedimientos: RAG sobre documentación oficial;
  \item preguntas sobre casos similares: recuperación de reclamos y transcripts históricos, con cautela;
  \item priorización de fraude, escalamiento o riesgo: inferencia estadística posterior.
\end{itemize}

El veredicto es, por tanto, \textbf{RAG híbrido más herramientas de consulta}, no RAG puro sobre las tablas ni inferencia estadística como primer mecanismo de respuesta.

\section{Criterio de evaluación}

Una fuente es apropiada para RAG cuando contiene texto suficientemente diverso, informativo, estable y autorizado para responder preguntas. Una fuente estructurada es más apropiada como herramienta cuando representa el estado actual de una entidad y la respuesta depende de filtros exactos, permisos o cálculos.

Se evaluaron tres conjuntos principales:

\begin{enumerate}
  \item \texttt{transactions}: operaciones financieras y señales transaccionales;
  \item \texttt{call\_transcripts}: conversaciones y anotaciones de soporte;
  \item \texttt{complaints}: reclamos, disputas, resoluciones y estados.
\end{enumerate}

La evaluación se hizo mediante conteos, cardinalidad textual, disponibilidad de vínculos y cobertura de campos, sin extraer ni exponer datos personales.

\section{Transacciones y operaciones}

La tabla \texttt{transactions} contiene 4.43 millones de registros. Para tarjetas de crédito hay 1\,108\,285 transacciones, con la siguiente distribución:

\begin{table}[ht]
\centering
\caption{Estados observados en transacciones de tarjetas de crédito}
\begin{tabular}{lrr}
\toprule
Estado & Registros & Proporción aproximada \\
\midrule
Aprobada & 1\,019\,373 & 92.0\% \\
Rechazada & 55\,650 & 5.0\% \\
Pendiente & 21\,952 & 2.0\% \\
Reversada & 11\,310 & 1.0\% \\
\bottomrule
\end{tabular}
\end{table}

También existen campos con valor para respuestas operativas:

\begin{itemize}
  \item \texttt{transaction\_status}: cuatro valores y sin nulos;
  \item \texttt{response\_code}: cinco valores distintos, aunque con aproximadamente 221 mil nulos en el total;
  \item \texttt{fraud\_score}: más de 5 mil valores distintos;
  \item \texttt{is\_fraud}, canal, monto, moneda y ubicación;
  \item relación con \texttt{customer\_id} y \texttt{product\_id}.
\end{itemize}

Estas columnas permiten responder preguntas como:

\begin{lstlisting}
Cual fue el estado de mi operacion?
Que operaciones fueron rechazadas?
Cual fue la ultima compra de mi tarjeta?
La operacion fue marcada como fraude?
Que codigo de respuesta tuvo la transaccion?
\end{lstlisting}

Para este tipo de consultas, un vector store puede introducir ambigüedad. La respuesta debe producirse con una consulta parametrizada y autorizada:

\[
\text{Pregunta} \rightarrow \text{identificación autorizada} \rightarrow \text{SQL} \rightarrow \text{respuesta con evidencia}
\]

Las transacciones pueden alimentar posteriormente modelos de fraude o priorización, pero no es necesario aplicar inferencia estadística para responder su estado factual.

\section{Transcripts de soporte}

La tabla \texttt{call\_transcripts} tiene 171\,321 transcripts. Aunque el volumen parece suficiente, la diversidad efectiva es baja:

\begin{table}[ht]
\centering
\caption{Calidad del corpus de transcripts}
\begin{tabularx}{\textwidth}{>{\raggedright\arraybackslash}p{0.55\textwidth} r}
\toprule
Métrica & Valor \\
\midrule
Transcripts totales & 171\,321 \\
Textos completos distintos & 546 \\
Longitud promedio & 278 caracteres \\
Mediana de longitud & 231 caracteres \\
Percentil 90 de longitud & 400 caracteres \\
Intenciones distintas & 1 \\
Temas distintos & 6 \\
Idiomas detectados & 1 \\
\bottomrule
\end{tabularx}
\end{table}

Para interacciones explícitamente relacionadas con tarjetas solo aparecen 322 transcripts, con 105 textos distintos y una única intención dominante: \texttt{consulta\_general}. Esto sugiere que el corpus es repetitivo o sintético y no representa suficientemente la variedad de problemas reales de soporte.

Los transcripts sí pueden servir para:

\begin{itemize}
  \item estudiar cómo formulan preguntas los clientes;
  \item construir ejemplos de prompts;
  \item crear un clasificador preliminar;
  \item recuperar conversaciones parecidas como contexto auxiliar.
\end{itemize}

No deberían usarse como fuente autorizada para afirmar políticas, plazos o acciones permitidas. Una conversación histórica no necesariamente representa la regla vigente del banco.

\section{Reclamos y disputas de tarjetas}

Hay 15\,486 reclamos vinculados a productos de tarjeta. El problema principal no es el número de filas, sino la escasa diversidad y trazabilidad textual:

\begin{table}[ht]
\centering
\caption{Cobertura de casos de disputa de tarjetas}
\begin{tabularx}{\textwidth}{>{\raggedright\arraybackslash}p{0.67\textwidth} r}
\toprule
Indicador & Casos \\
\midrule
Casos asociados a tarjetas & 15\,486 \\
Descripciones distintas & 5 \\
Resoluciones distintas & 5 \\
Casos con texto de resolución & 3\,534 \\
Casos con satisfacción registrada & 561 \\
Casos con monto reclamado & 5\,036 \\
Casos con compensación & 1\,053 \\
Casos con fecha de cierre & 565 \\
Casos con \texttt{origin\_interaction\_id} & 0 \\
\bottomrule
\end{tabularx}
\end{table}

Además, la tabla de reclamos no contiene un \texttt{transaction\_id}. Por tanto, no es posible enlazar directamente una disputa con la operación concreta que el cliente cuestiona. El vínculo existente es principalmente con el producto de tarjeta y con campos de categoría, monto, estado, prioridad y fechas.

Esto permite responder mediante SQL:

\begin{itemize}
  \item estado actual del reclamo;
  \item prioridad y categoría;
  \item canal de recepción;
  \item monto reclamado y compensación;
  \item fecha de resolución o cierre, cuando existe;
  \item incumplimiento de SLA.
\end{itemize}

Pero limita las respuestas narrativas como:

\begin{itemize}
  \item “¿Por qué se rechazó mi disputa?”;
  \item “¿Qué ocurrió exactamente con mi compra?”;
  \item “¿Qué documentos debo presentar?”;
  \item “¿Cuál es la política para una operación no reconocida?”.
\end{itemize}

Las cinco descripciones y cinco resoluciones distintas sugieren que los textos son plantillas o datos sintéticos muy repetitivos. Pueden utilizarse para prototipos, pero no ofrecen un corpus robusto para recuperación semántica.

\section{Evaluación por tipo de pregunta}

\begin{table}[ht]
\centering
\small
\caption{Mecanismo recomendado por tipo de pregunta}
\begin{tabularx}{\textwidth}{>{\raggedright\arraybackslash}p{0.29\textwidth} >{\raggedright\arraybackslash}p{0.25\textwidth} X}
\toprule
Tipo de pregunta & Mecanismo & Viabilidad \\
\midrule
Estado de tarjeta, saldo, límite o mora & SQL/herramienta & Alta; depende de autorización y filtros exactos. \\
Estado de una operación & SQL/herramienta & Alta; existe estado y contexto transaccional. \\
Motivo de rechazo & SQL + reglas & Media-alta; depende de la cobertura de \texttt{response\_code}. \\
Señal de fraude & SQL + modelo posterior & Media-alta; existe evidencia estructurada, pero no siempre explicación. \\
Estado de un reclamo & SQL/herramienta & Media-alta; hay estado, prioridad y fechas. \\
Política de disputa & RAG documental & Baja con BigQuery actual; falta documentación oficial. \\
Caso histórico similar & RAG de casos & Baja-media; el corpus es pequeño y repetitivo. \\
Acción permitida por el agente & RAG + reglas & Baja con los datos actuales; faltan políticas y permisos. \\
\bottomrule
\end{tabularx}
\end{table}

\section{Arquitectura recomendada}

La solución debería separar las responsabilidades:

\begin{lstlisting}
Pregunta del cliente
        |
        +-- Pide estado o dato actual?
        |       |
        |       +-- Consulta autorizada a BigQuery
        |
        +-- Pide politica o procedimiento?
        |       |
        |       +-- RAG sobre documentacion oficial
        |
        +-- Pide comparacion con casos anteriores?
                |
                +-- Recuperacion de casos historicos
\end{lstlisting}

Las herramientas estructuradas podrían exponer interfaces como:

\begin{lstlisting}
get_customer_cards(customer_id)
get_card_status(product_id)
get_recent_transactions(customer_id, product_id, date_range)
get_transaction_details(transaction_id)
get_card_disputes(customer_id, product_id)
get_dispute_status(complaint_id)
\end{lstlisting}

El RAG documental debería contener políticas versionadas sobre bloqueo, reemplazo, fraude, disputas, límites, requisitos, SLA y diferencias por país. Ese contenido no está suficientemente presente en BigQuery.

\section{Dónde encaja la inferencia estadística}

La inferencia no debe ser el primer mecanismo para responder hechos actuales. Puede incorporarse como una capa de apoyo para:

\begin{itemize}
  \item priorizar casos de fraude;
  \item identificar disputas con necesidad probable de escalamiento;
  \item detectar patrones anómalos;
  \item recomendar una categoría de problema;
  \item ordenar casos para atención humana.
\end{itemize}

Una formulación posterior podría estimar:

\[
P(\text{Fraud}\mid\text{channel},\text{amount},\text{fraud score},\text{response code},\text{location},\ldots)
\]

Pero el modelo debería devolver evidencia y confianza, no reemplazar la fuente transaccional ni la política oficial.

\section{Riesgos y controles}

\begin{itemize}
  \item \textbf{Información sensible:} no indexar directamente números de producto, documentos, teléfonos o correos; aplicar autorización por cliente.
  \item \textbf{Datos actuales frente a históricos:} separar claramente la consulta del estado actual de la recuperación de casos antiguos.
  \item \textbf{Corpus repetitivo:} deduplicar transcripts y reclamos antes de cualquier indexación.
  \item \textbf{Falta de vínculo transaccional:} agregar \texttt{transaction\_id} a disputas y enlazar las interacciones originales.
  \item \textbf{Políticas ausentes:} no inferir requisitos o decisiones operativas a partir de ejemplos históricos.
  \item \textbf{Trazabilidad:} cada respuesta debe indicar si proviene de una fila actual, una política documental o un caso histórico.
\end{itemize}

\section{Recomendaciones de implementación}

\begin{enumerate}
  \item Construir primero un asistente de consulta para tarjetas, transacciones y disputas.
  \item Crear una base RAG separada con documentación oficial y versionada.
  \item Usar transcripts y reclamos únicamente como ejemplos históricos filtrables.
  \item Normalizar y deduplicar los textos antes de generar embeddings.
  \item Crear una relación explícita entre conversación, tarjeta, transacción, reclamo y resultado.
  \item Añadir una taxonomía de intenciones y códigos de resolución específicos de tarjetas.
  \item Introducir inferencia estadística solo para priorización, riesgo y escalamiento.
\end{enumerate}

\section{Conclusión}

BigQuery sí puede fortalecer una solución de pregunta--respuesta para Dispute Transaction Support, pero no como un RAG puro sobre todas las tablas.

La arquitectura más defendible es:

\begin{quote}
\textbf{BigQuery como fuente estructurada y autorizada para datos actuales de tarjetas, operaciones y disputas; RAG sobre políticas oficiales; y transcripts/reclamos como evidencia histórica secundaria.}
\end{quote}

Este enfoque permite comenzar con respuestas verificables sin depender de inferencia estadística general. La inferencia puede añadirse posteriormente para detectar riesgo, priorizar casos y recomendar escalamiento.

\end{document}
